\documentclass[10pt,letterpaper]{article}
\usepackage[margin=1in]{geometry}
\usepackage[T1]{fontenc}
\usepackage{amsmath,amssymb,amsthm,booktabs,longtable,array,microtype}
\usepackage[round,authoryear]{natbib}
\usepackage{xurl}
\usepackage{placeins}
\usepackage[hidelinks]{hyperref}
\hypersetup{pdftitle={Supporting methods: When More Search Can Worsen Deployment: Best-of-N Selection on an Incomplete Proxy},pdfauthor={Dev D. Goyal}}
\newcommand{\E}{\mathbb E}
\newcommand{\Prb}{\mathbb P}
\newcommand{\ind}{\mathbf 1}
\title{Supporting methods for\\When More Search Can Worsen Deployment:\\Best-of-$N$ Selection on an Incomplete Proxy}
\author{Dev D. Goyal}
\date{}
\begin{document}
\maketitle
\noindent This supplement explains how we construct the trading candidates, simulate their outcomes, and calculate the effect of search. Sections~\ref{sec:data}--\ref{sec:execution} give the data, forecasts, and execution model. Section~\ref{sec:uncertainty} gives the selection rules, resampling procedures, and sensitivity results. Section~\ref{sec:status} states the study's retrospective limits. The source package includes aggregate results. It does not redistribute the vendor price records.

\section{Finance data and feature construction}\label{sec:data}
The daily price and volume records came from the Yahoo Finance chart endpoint, \url{https://query1.finance.yahoo.com/v8/finance/chart/}. The retained data end on 24 July 2026. The download starts on 1 January 1998. The analysis panel starts on 4 January 1999. This panel contains daily records for a fixed set of 258 US equity tickers. The download script starts from a fixed list of 347 identifiers. The inspected scripts do not identify an external source for that initial list. The original panel retains stocks with more than 90\% nonmissing adjusted-price coverage from 1 January 2009 through the retained endpoint. The extended panel preserves that original ticker set. The exact 258 identifiers, in panel order, are in \texttt{data/finance\_universe.csv}. This is a survivor-filtered universe. It does not reconstruct all historical index members or delisting returns.

If a ticker has multiple records for one date, we retain the last record. Let $P_{it}$ denote adjusted close and $D_{it}$ denote provider close times reported volume. We fill missing adjusted closes with the last available value for at most three panel rows. Simple returns are
\[
 r_{it}=\operatorname{clip}(P_{it}/P_{i,t-1}-1,-0.5,0.5),\qquad
 V_{it}=\operatorname{median}_{21}(D_i)_t.
\]
The percentage-change calculation does not fill any further missing values. The median requires at least 10 observations. We use $V$ to estimate daily dollar volume from the trailing observations. It is a median, although the implementation calls it ADV, for average daily volume. All row counts below refer to the common panel calendar.

Define $\ell_{it}=\log P_{it}$ and $m_t$ as the mean available return across retained tickers. We transform each base feature $x$ into a daily percentile rank minus one half:
\[
 z_{it}(x)=\operatorname{rank}_{t}(x_{it})/n_t(x)-\tfrac12 .
\]
Here $n_t(x)$ is the number of nonmissing feature values on date $t$. Ties receive average ranks. Missing values remain missing. Table~\ref{tab:features} gives the 19 raw formulas. Every named rolling window requires its full number of nonmissing observations, except $V$, which retains its 10-of-21 rule. Residual volatility requires the most history: 314 rows after a ticker's first record.

\begin{table}[t]
\centering
\caption{Raw base features in stored order. Each is transformed to a daily percentile rank minus one half. Rolling means, sample standard deviations, maxima, covariances, and variances use full named windows. The symbol $\beta_{it}$ in the residual formula denotes the raw rolling beta, before its rank transform.}
\label{tab:features}
\small
\begin{tabular}{@{}ll@{}}
\toprule
Name & Raw formula through close $t$\\
\midrule
\texttt{mom21} & $\ell_{it}-\ell_{i,t-21}$\\
\texttt{mom63} & $\ell_{it}-\ell_{i,t-63}$\\
\texttt{mom126} & $\ell_{i,t-21}-\ell_{i,t-126}$\\
\texttt{mom252} & $\ell_{i,t-21}-\ell_{i,t-252}$\\
\texttt{rev1} & $r_{it}$\\
\texttt{rev5} & $\ell_{it}-\ell_{i,t-5}$\\
\texttt{vol21} & $s_{21,t}(r_i)$\\
\texttt{vol63} & $s_{63,t}(r_i)$\\
\texttt{ladv} & $\log\max(V_{it},1)$\\
\texttt{dvolz} & $\log\max(V_{it},1)-\overline{\log\max(V_i,1)}_{63,t}$\\
\texttt{amihud} & $10^9\,\overline{|r_i|/\max(V_i,1)}_{21,t}$\\
\texttt{p52w} & $P_{it}/\max_{252,t}(P_i)$\\
\texttt{beta} & $\beta_{it}=\operatorname{Cov}_{252,t}(r_i,m)/\operatorname{Var}_{252,t}(m)$\\
\texttt{ivol} & $s_{63,t}(r_i-\beta_i m)$\\
\texttt{skew63} & $\operatorname{Skew}_{63,t}(r_i)$\\
\texttt{maxr21} & $\max_{21,t}(r_i)$\\
\texttt{ma50} & $P_{it}/\overline P_{i,50,t}$\\
\texttt{ma200} & $\overline P_{i,50,t}/\overline P_{i,200,t}$\\
\texttt{accel} & $(\ell_{it}-\ell_{i,t-21})-(\ell_{it}-\ell_{i,t-63})/3$\\
\bottomrule
\end{tabular}
\end{table}


We construct the other ten features from the stored ranks. They are not ranked again. Five are $z(x)^2-1/12$ for $x\in\{\texttt{mom252},\texttt{rev1},\texttt{vol21},\texttt{p52w},\texttt{ladv}\}$. Five are products for the ordered pairs
\[
\begin{gathered}
 (\texttt{mom252},\texttt{vol21}),\quad(\texttt{rev1},\texttt{vol21}),\quad
 (\texttt{mom21},\texttt{ladv}),\\
 (\texttt{p52w},\texttt{amihud}),\quad(\texttt{mom252},\texttt{rev1}).
\end{gathered}
\]
The feature order is the table order, followed by the five squares in the order above, then the five products in the order above. This order determines the seeded random feature sets.

We fix the stock universe before constructing features. We then construct and rank the features before checking daily trading eligibility. Thus $m_t$ and a daily feature rank can include a retained ticker that later fails that day's trading rules. A ticker on a given date is valid only if all 29 features and its return are present, $V_{it}>2\times10^6$, and $P_{it}>5$. The last test uses adjusted close. It is a retrospective adjusted-price screen, not a verified historical nominal USD 5 screen. We store missing numeric entries as zero and use a separate validity mask to record eligibility. We multiply features by that mask before estimation and scoring. No industry adjustment, feature winsorization, or feature z-score is applied. The return clip occurs before feature construction.

\section{Configurations and pooled ridge forecasts}\label{sec:fit}
Each candidate configuration combines six choices. Taking every combination gives
$4\times4\times4\times25\times5\times3=24{,}000$ entries for the feature sets constructed in each pool.
\begin{center}
\begin{tabular}{@{}ll@{}}
\toprule
Axis & Values\\
\midrule
Forecast horizon $h$ & $1,5,10,21$ rows\\
Training window $W$ & $252,504,756,1260$ rows\\
Ridge parameter $\lambda$ & $10^{-3},10^{-2},10^{-1},1$\\
Feature set $S$ & 11 named sets and 14 seeded random sets\\
Score half-life $H$ & $0,2,5,10,21$ rows\\
Portfolio transform & \texttt{rank}, \texttt{q20}, \texttt{q10}\\
\bottomrule
\end{tabular}
\end{center}
Pools A and D contain 500 configurations each. Pools B and C contain 1,500 each. The NumPy default generator uses seeds 11, 202, 303, and 404 for pools A through D, respectively. The generator first constructs all 14 random feature sets. Each set draws its size uniformly from the integers 6 through 19, then draws that many feature indices without replacement and sorts them. The generator draws each set separately. It can draw duplicate sets, and no rule limits overlap between sets.

Let $\mathcal L$ contain all 19 base features and $\mathcal N$ contain all ten nonlinear features. We also define four feature groups:
\begin{align*}
\mathcal M={}&\{\texttt{mom21},\texttt{mom63},\texttt{mom126},\texttt{mom252},
\texttt{accel},\texttt{ma50},\texttt{ma200},\texttt{p52w}\},\\
\mathcal R={}&\{\texttt{rev1},\texttt{rev5},\texttt{maxr21},\texttt{skew63}\},\\
\mathcal K={}&\{\texttt{vol21},\texttt{vol63},\texttt{ivol},\texttt{beta}\},\\
\mathcal Q={}&\{\texttt{ladv},\texttt{dvolz},\texttt{amihud}\}.
\end{align*}
In insertion order, the named feature sets are $\mathcal L$, $\mathcal L\cup\mathcal N$, $\mathcal M\cup\mathcal R\cup\mathcal N$, $\mathcal M$, $\mathcal R$, $\mathcal K$, $\mathcal Q$, $\mathcal M\cup\mathcal R$, $\mathcal M\cup\mathcal K$, $\mathcal R\cup\mathcal Q$, and $\mathcal L\setminus\mathcal Q$. The corresponding names are \texttt{lin\_all}, \texttt{all}, \texttt{nl\_only}, \texttt{mom}, \texttt{rev}, \texttt{risk}, \texttt{liq}, \texttt{mom+rev}, \texttt{mom+risk}, \texttt{rev+liq}, and \texttt{no\_liq}. Despite its name, \texttt{nl\_only} includes linear momentum and reversal features. The 14 random sets follow these named sets.

After constructing the feature sets, the same generator samples one value from each row of the table, in table order. We reject any configuration already in the pool. Sampling continues until the pool has the required number of configurations. Named sets are common across pools. Random sets depend on the pool seed. Thus the four pools are configuration draws on one shared market history, not independent market replications.

The model predicts a label for each stock at decision row (origin) $s$. This label is the forward log return over $h$ rows, starting at the next fill close:
\[
 g_{is}^{(h)}=\sum_{u=s+2}^{s+1+h}\log(1+r_{iu}).
\]
We include the label when the trading-validity masks at $s+1$ and $s+1+h$ are both true. We do not test the masks between those endpoints. Since missing returns are stored as zero, an intermediate missing return contributes zero when both endpoint tests pass. On each origin date, we subtract the mean of the included labels and divide by their population standard deviation. If at most 20 names qualify, or that standard deviation is zero, all labels on that origin are set to zero. Write the standardized labels as $y_{is}$. Labels overlap across origins. We separately set the feature vector at origin $s$ to zero if $s$ is invalid.

We refit each model every 21 rows, starting on 1 July 2004, the first scoring row. The coefficients remain fixed between refits. At refit row $t$, set $e=t-h-1$. The included origins are exactly $e-W+1,\ldots,e$. The last label ends at close $t$. Each refit requires a full training window. We exclude no additional dates between training and scoring. The counts of valid label endpoints exceed 20 on every included training date in the retained US calculation. The smallest count is 129 across the 265 refits.

Let $x_{is,S}$ be the selected feature vector after its origin validity mask, and let $a_{is}^{(h)}$ indicate the two valid label endpoints. Define
\[
 M_t=\sum_{s=e-W+1}^{e}\sum_i a_{is}^{(h)}x_{is,S}x_{is,S}^{\mathsf T},\qquad
 b_t=\sum_{s=e-W+1}^{e}\sum_i a_{is}^{(h)}x_{is,S}y_{is}.
\]
The fitted coefficient solves
\[
 \left[M_t+\lambda\left\{\operatorname{tr}(M_t)/|S|+10^{-12}\right\}I\right]
 \widehat\beta_t=b_t.
\]
There is no intercept, observation weighting, or training-window feature standardizer. We scale the ridge penalty by the mean diagonal of the Gram matrix. We do not separately normalize it by the sample count. Cumulative Gram matrices use double precision. Each refit uses a direct linear solve. If the linear system is singular, the solver returns zero coefficients. Stored features and scores use single precision.

The raw forecast is $f_{it}=x_{it,S}^{\mathsf T}\widehat\beta_t$ and is missing on invalid rows. For $H>0$, forecasts use the pandas exponentially weighted mean with $\alpha=1-0.5^{1/H}$, \texttt{min\_periods=1}, \texttt{ignore\_na=True}, and the default adjusted-weight convention. $H=0$ bypasses smoothing. We smooth the forecasts before constructing the portfolio. We apply the validity mask again before ranking the daily scores.

\section{Target holdings, fills, and daily outcomes}\label{sec:execution}
Let $\rho_{it}$ be the percentile rank of a valid smoothed forecast. Define the preliminary score by
\[
 a_{it}=\begin{cases}
 \rho_{it}-0.5,&\texttt{rank},\\
 \ind\{\rho_{it}\geq0.8\}-\ind\{\rho_{it}\leq0.2\},&\texttt{q20},\\
 \ind\{\rho_{it}\geq0.9\}-\ind\{\rho_{it}\leq0.1\},&\texttt{q10}.
 \end{cases}
\]
Set missing scores to zero. Subtract the mean score over valid names and divide by the sum of absolute demeaned scores, floored at $10^{-9}$. This gives target holdings $w_{it}$ with zero net exposure and unit gross exposure when the demeaned scores have total absolute value at least as large as the floor. If all scores are zero, all holdings remain zero. The long and short legs are not normalized separately. Under the quantile transforms, unequal tail counts can give small nonzero weights to middle-ranked names after demeaning. There is no maximum name weight, sector constraint, beta constraint, covariance model, or volatility target.

We express each holding as its dollar value divided by fixed target capital $A$. These holdings are not fractions of current wealth. A decision uses information through close $t$. The trade fills at the next close, $t+1$. Its post-fill holdings earn the return $r_{i,t+2}$ from close $t+1$ to close $t+2$. The final two decision rows without a subsequent realized return are excluded. The last scored decision date is 22 July 2026.

Holdings start at zero on 1 July 2004. Before each trade, we update the value of existing holdings for the latest price change. Thus $v_{it}=\widehat w_{i,t-1}(1+r_{i,t+1})$. Let $d_{it}=w_{it}-v_{it}$ be the gap to the new target and $\overline V_{it}=\max(V_{it},10^6)$ the execution volume input. Execution is
\[
 e_{it}=\operatorname{clip}\left(d_{it},-\kappa \overline V_{it}/A,\kappa \overline V_{it}/A\right),
 \qquad \widehat w_{it}=v_{it}+e_{it},\qquad \kappa=0.10.
\]
The uncapped gross state uses an infinite cap and no charge. Other states use $A\in\{25,100,500,2000\}$ million USD. We do not carry unfilled orders forward. At the next decision, we compare the holdings with the new target. An ineligible stock can remain in the portfolio when the trading cap limits its sale or purchase to close the position. We do not hedge any remaining net exposure.

The cost model uses information available through the signal close $t$. The execution model sets a lower limit of USD $10^6$ for $V$. Its volatility $\sigma_{it}$ is the sample standard deviation of valid returns over 21 rows, with at least ten observations. Missing volatility is set to 0.02. Volatility is then clipped to $[0.005,0.15]$. Let
\[
 G_t=\max\left(\sum_i|\widehat w_{it}|,10^{-9}\right),\qquad
 p_{it}=|e_{it}|A/\overline V_{it}.
\]
The stored daily components are
\begin{align*}
 R_t^{g}(A)&=G_t^{-1}\sum_i\widehat w_{it}r_{i,t+2},\\
 T_t(A)&=G_t^{-1}\sum_i|e_{it}|,\qquad
 I_t(A)=G_t^{-1}\sum_i\sigma_{it}|e_{it}|\sqrt{p_{it}},\\
 R_t^{n}(A)&=R_t^{g}(A)-sT_t(A)-\eta I_t(A),
 \quad s=1.5\times10^{-4},\quad\eta=0.5.
\end{align*}
Each series divides its daily amount by the gross exposure actually held. Spread and impact change charges but do not change the executed holdings at a fixed $A$ and $\kappa$. The square-root impact function is an assumed scenario consistent with common impact models \citep{almgren2005direct,toth2011anomalous}. It is not calibrated to observed fills in this study.

Partial fills can leave nonzero net exposure. Financing, stock-borrow charges, margin requirements, and wealth-based position reduction are absent. Diagnostic cash pays for trades and the default modeled charges. Cash and diagnostic wealth do not change the targets, trading caps, or return normalization. Negative diagnostic wealth does not terminate a path. Therefore these series do not give the return or excess return of a funded investor account.

\section{Selection curves, cost-aware rules, and resampling}\label{sec:uncertainty}
The primary selection window is $[\text{4 January 2010},\text{2 January 2019})$ as measured by decision dates. Evaluation starts on 2 January 2019 and ends at the last scored row. Returns from the final selection decisions can be realized after the first evaluation decision, so this boundary is a retrospective outcome partition, not a ranking frozen prospectively at that decision. Models and holdings continue through these boundaries. There is no reset to cash at the evaluation boundary. Each configuration uses the same scored dates.

For a daily series $R$, its reported Sharpe is $\sqrt{252}\,\overline R/\operatorname{sd}(R)$ with population standard deviation. A zero standard deviation gives zero Sharpe. No risk-free return is subtracted. The baseline ranks candidates by selection-window uncapped gross Sharpe, with larger values preferred. For each fixed ranking, exact without-replacement rank weights give expected selected evaluation Sharpe at $N\in\{1,2,4,8,16,32\}$. We fit a line to these means against $\log_2N$ by ordinary least squares, with an intercept. Its slope is $\Lambda$. This slope describes the six-point grid. It does not imply an equal effect at every doubling.

Rank a pool of $m$ candidates from best to worst by the selection score. The implementation sorts scores in ascending order with a stable sort. The later stored configuration wins an exact tie. The rank convention here reverses that order, so rank 1 is best. Under uniform sampling without replacement, the probability that rank $k$ wins is
\[
 p_N(k)=\frac{\binom{m-k}{N-1}}{\binom{m}{N}},\qquad 1\leq k\leq m.
\]
A binomial coefficient is zero when its lower argument exceeds its upper argument. A candidate wins if it is sampled and all other sampled candidates have worse ranks. For any fixed candidate quantity $v_k$, the selected mean is $\sum_k p_N(k)v_k$. This standard order-statistic calculation gives the mean without repeatedly simulating searches \citep{oneill2023order}. The code evaluates the weights with log-gamma functions and normalizes their sum. The identity is exact. The numerical calculation uses floating-point arithmetic.

The same weights give expected turnover and the probability that the selected candidate has positive evaluation Sharpe. To find a selected median, order the outcomes by value and add their winner probabilities. The median is the first outcome value where that sum reaches one half. These are summaries of candidate outcomes. In particular, expected candidate Sharpe is not the Sharpe of an averaged portfolio. The primary grid searches at most 10\% of each pool. This is a study-design restriction, not a restriction of the rank identity.


At each positive capital, we evaluate four paths in this order: uncapped gross, capped gross, capped with spread, and capped with spread and impact. Write their fitted slopes as $\Lambda_0,\Lambda_c,\Lambda_s,\Lambda_n$. The total gap slope is $\Gamma=\Lambda_0-\Lambda_n$. Its capacity-related component is
\[
 \Gamma_{\rm cap}=(\Lambda_0-\Lambda_c)+(\Lambda_s-\Lambda_n).
\]
These differences separate the spread step from the capping and impact steps, in the stated order. This is not a unique decomposition of cash costs. The signed Sharpe gap can be negative. Positive and negative local rank contributions can cancel.
Here capacity means the effect of a fixed target capital on the participation limit and modeled price impact. It is not a measured maximum investment size. The adjacent slope differences add to the total because the intermediate slopes cancel. A positive gap slope means that the difference between gross and net Sharpe grows. It does not by itself mean that net Sharpe falls.

We also test whether a different selection score improves the outcome under the same best-of-$N$ rule. This comparison is exploratory. For assumed capital $A'\in\{25,100,500\}$ million USD, one score is selection-period net Sharpe at $A'$. The calibrated turnover score is
\[
 c_i(A')=S_i^{g,\mathrm{sel}}-\gamma(A')\,\overline T_i^{\mathrm{sel}}(A'),
 \qquad
 \gamma(A')=
 \frac{m^{-1}\sum_i\{S_i^{g,\mathrm{sel}}-S_i^{n,\mathrm{sel}}(A')\}}
 {\max\{m^{-1}\sum_i\overline T_i^{\mathrm{sel}}(A'),10^{-12}\}}.
\]
We calibrate this coefficient separately for each pool and assumed capital. It uses selection-window gross Sharpe, modeled net Sharpe, and executed turnover. Thus this penalty requires the execution model. It is not a model-free turnover correction. We average expected net Sharpe at $N=32$ over four pools, three assumed capitals, and three evaluation capitals, for 36 equally weighted cells. The cells share market paths and candidates. Their mean is not an estimate from 36 independent deployments.

A uniform bound on score error would limit the loss from choosing the wrong candidate. Suppose $|\widehat u_i-u_i|\leq\epsilon$ for every candidate in a subset. If $a$ maximizes $u$ and $b$ maximizes $\widehat u$, then
\[
 u_a-u_b=(u_a-\widehat u_a)+(\widehat u_a-\widehat u_b)
                  +(\widehat u_b-u_b)\leq2\epsilon.
\]
The study does not measure this uniform bound for Sharpe. A bound on dollar-cost error alone would not establish it.

Both finance resampling analyses use 1,000 stationary block bootstrap replicates. We report the 2.5th and 97.5th percentiles. The percentile interval contains the middle 95\% of the replicate values. Each index vector starts with a day drawn uniformly from the window. Each next position either continues to the next day or restarts at a uniformly drawn day with probability $1/L$. A continuing block wraps from the last day to the first day when needed. The resulting list of day indices is the stationary index vector. The primary mean block length is $L=21$. Selection and evaluation have separate index vectors. Within a window, all candidates use the same vector. The bootstrap resamples stored daily outcome components. It does not reconstruct prices, refit the ridge models, regenerate portfolios, or rerun fills.

For the pool-B capacity-component interval, joint resampling recomputes selection Sharpe, rankings, rank weights, and slopes. Evaluation-only resampling holds the selection ranking fixed. The base seed is 20260807. For each mode and block length, the implementation initializes NumPy's default generator with $20260807+L$, hence 20260828 at $L=21$. Within each joint replicate, evaluation indices are drawn before selection indices. The block-length sensitivities use $L\in\{5,10,21,63\}$.

The comparison of cost-aware rules uses seed 20260825 and $L=21$. It draws selection indices before evaluation indices. The same two index vectors are shared across all pools, candidates, and selectors. Each replicate recomputes selection Sharpes, turnover, $\gamma$, rankings, and paired differences in the 36-cell mean. The intervals are descriptive resampling intervals conditional on the stored models, execution assumptions, and observed history. They are not validated population confidence guarantees. The capacity component's joint interval includes zero. Mean net Sharpe remains negative under both cost-aware scores.

\subsection{Complete exploratory selector comparison}\label{sec:all-selectors}
Table~\ref{tab:all-selectors} lists all seven tested rules. The two simulator-free penalty coefficients are variants of one rule. All values come from retained aggregate records for the simulation that updates holdings for price changes. The accompanying \texttt{data/finance\_additional\_checks.json} preserves those records, their source hashes, and the arithmetic used for the summaries. Only the net-score and calibrated-turnover rules have paired gain intervals. These comparisons are exploratory.

\begin{table}[htbp]
\centering
\begin{tabular}{@{}lrr@{}}
\toprule
Selection rule & Mean net Sharpe & Gain over gross selection\\
\midrule
Gross Sharpe & $-0.501$ & $0.000$\\
Turnover decile screen & $-0.298$ & $+0.203$\\
Simulator-free penalty, $\eta_{\rm ref}=0.5$ & $-0.169$ & $+0.332$\\
Simulator-free penalty, $\eta_{\rm ref}=0.11$ & $-0.229$ & $+0.272$\\
Calibrated turnover penalty & $-0.048$ & $+0.453$\\
Net Sharpe at assumed capital & $-0.071$ & $+0.430$\\
Deflated net Sharpe at assumed capital & $-0.071$ & $+0.430$\\
Ex-post evaluation oracle & $+0.220$ & $+0.721$\\
\bottomrule
\end{tabular}
\caption{The full retrospective comparison at $N=32$ uses the same 36 equally weighted combinations of pool, assumed capital, and evaluation capital. All non-oracle rules have negative mean net Sharpe. The oracle uses the evaluation outcomes and is unavailable at selection. The simulator-free penalty uses full-panel universe summaries. These summaries can include evaluation dates. Thus the table does not compare seven selection-window-only policies.}
\label{tab:all-selectors}
\end{table}
\FloatBarrier

Gross selection ranks the full pool by selection-period uncapped Sharpe. The turnover decile screen first removes candidates above the 90th percentile of selection-period executed turnover at the assumed capital. It draws $N$ candidates from the retained pool and then applies gross selection. This changes the candidate distribution relative to the full-pool scoring rules. This rule needs the fill model, even though its score uses no cost charge. The calibrated turnover penalty uses the coefficient defined above. The net-score rule uses selection-period net Sharpe at the assumed capital.

The simulator-free penalty ranks candidates by gross Sharpe minus $\gamma_{\rm sf}$ times desired turnover. Here desired turnover is turnover on the uncapped portfolio path. Its coefficient is
\[
 \gamma_{\rm sf}(A')=
 \frac{\sqrt{252}\{s+\eta_{\rm ref}\sigma_{\rm name}
           \sqrt{\tau A'/(n_{\rm live}V_{\rm med})}\}}
      {\sigma_{\rm book}}.
\]
Here $\tau$ is the median desired turnover across the pool, measured in the selection period. The denominator is the median daily volatility of the uncapped portfolios in that period. The other inputs summarize the full stored panel. They are median stock volatility $\sigma_{\rm name}=0.01699$, median daily dollar volume $V_{\rm med}=151.75$ million USD, and the mean number of valid stocks $n_{\rm live}=218.86$. These summaries can include evaluation dates. Thus the stored rule is retrospective. Its selection score needs no simulated fill path, but its reported net outcome still needs the simulator. The lower coefficient matches the temporary-impact form $0.142\sigma p^{0.6}$ to $\eta\sigma p^{1/2}$ at $p=0.10$ \citep{almgren2005direct}. The matching coefficient is $\eta(p)=0.142p^{0.1}$: about $0.113$ at $p=0.10$ and $0.071$ at $p=0.001$. The sensitivity uses the rounded value $0.11$. Both reference coefficients, $0.5$ and $0.11$, are retained in Table~\ref{tab:all-selectors}.

The deflated-net rule applies the Deflated Sharpe Ratio \citep{bailey2014deflated} to each selection-period net series. Let $z_i$ be its daily Sharpe, $s_z$ the sample standard deviation of $z_i$ across candidates, and $n$ the common observation count. Let $a_i$ and $b_i$ be the series' skewness and non-excess kurtosis. Non-excess kurtosis is 3 for a normal distribution. The trial count is 32. With $c_E=0.5772156649$ and standard normal distribution function $\Phi$, the implemented score is
\[
 z_* = s_z\{(1-c_E)\Phi^{-1}(1-1/32)+c_E\Phi^{-1}(1-1/(32e))\},
\]
\[
 \operatorname{DSR}_i=\Phi\!\left(
 \frac{(z_i-z_*)\sqrt{\max(n-1,1)}}
 {\sqrt{\max\{1-a_i z_i+(b_i-1)z_i^2/4,10^{-8}\}}}\right).
\]
We use this score to rank candidates. It can therefore change the selected candidate. Its mean selected net Sharpe is $-0.071231$, compared with $-0.071205$ for net selection. Equal values at three decimal places do not establish identical rankings or identical winners for every subset. The ex-post oracle ranks candidates by evaluation-period net Sharpe at the evaluation capital. It gives the best possible winner within each sampled subset for this recorded outcome. It is a diagnostic upper bound, not a usable selector.

\subsection{Reported sensitivity checks}\label{sec:finance-inventory}
Table~\ref{tab:finance-inventory} reports the other sensitivity checks, including results that limit the primary findings. Unless stated otherwise, the values use the same updates to holdings for price changes as the primary results. The data file preserves full aggregate rows, including cells outside the displayed capital range. The stored sign convention is $D=-\Gamma$ and $D_{\rm capacity}=-\Gamma_{\rm cap}$. Interval endpoints reverse order when their signs are changed. Gross slopes and selector gains retain their recorded signs.

\begin{longtable}{@{}p{0.25\linewidth}p{0.70\linewidth}@{}}
\caption{Retained finance checks and their scope. Counts describe dependent scenarios on shared data. They do not estimate success probabilities across deployments.}\label{tab:finance-inventory}\\
\toprule
Check & Retained result and interpretation\\
\midrule
\endfirsthead
\toprule
Check & Retained result and interpretation\\
\midrule
\endhead
Gross-slope uncertainty & Pool B has $\Lambda_{\rm gross}=0.0153$. At $L=21$, its joint interval is $[-0.0438,0.0678]$. The evaluation-only interval is $[-0.0240,0.0558]$. Both include zero. These resampling results do not establish that more search improves gross Sharpe.\\
Gap uncertainty & Joint and evaluation-only modes, four block lengths $5,10,21,63$, and the recorded capital levels are retained. The joint $L=21$ interval for $\Gamma_{\rm cap}(500\mathrm{M})$ is $[-0.0219,0.1209]$. It includes zero.\\
Search grids & Five grids anchored at $N=1$ give 64 eligible pool-by-grid-by-capital cells. All have positive $\Gamma_{\rm cap}$. Their 500M range is $[0.0172,0.0571]$. This count excludes the separate late-start check. Pool A on $\{8,16,32\}$ has $\Gamma_{\rm cap}(500\mathrm{M})=-0.0070$.\\
Execution factorial & The tested participation limits are $0.05,0.10,0.20$. Spreads are $0.75,1.5,3$ basis points. Impact coefficients are $0.25,0.5,1$. All 324 cells at 25M, 100M, and 500M have positive $\Gamma_{\rm cap}$. The 500M range is $[0.0069,0.1106]$. Across all three capitals the lower endpoint is $0.0036$.\\
Walk-forward periods & Six selection--evaluation pairs give 72 pool-by-period-by-capital cells. Of these, 65 have positive $\Gamma_{\rm cap}$. The 500M range is $[-0.0096,0.2418]$. Only 16 of 24 pool-by-period gross slopes are positive. Net selection at assumed 25M improves the level and slope in all 72 cells.\\
Cost misspecification & The 2,916 paired execution scenarios use assumed capital 25M and evaluation capital 500M. All have positive net-score gains over gross selection. The minimum is $0.1234$. These are scenario comparisons, not independent confirmations.\\
Impact coefficient & The temporary-impact mapping $\eta(p)=0.142p^{0.1}$ gives about $0.11$ at $p=0.10$ and $0.07$ at $p=0.001$. At $\eta=0.11$, all 12 pool-by-capital cells at 25M, 100M, and 500M have positive $\Gamma_{\rm cap}$. The 500M range is $[0.0052,0.0106]$. Other retained coefficients include $0.07,0.25,0.5,1$.\\
Impact exponent & Replacing $p^{1/2}$ by $p^{0.6}$ and matching charges at $p=0.10$ gives $\eta_{0.6}=0.62946$. The 500M $\Gamma_{\rm cap}$ range becomes $[0.0168,0.0424]$. Mean net-score gain falls from $0.4298$ to $0.3781$. The full-pool best candidate under each score is unchanged by the exponent change in all 12 pool-by-assumed-capital cases. This does not imply that every subset winner stays the same.\\
Exposure diagnostics & Under gross selection at $N=32$ and 500M, pool ranges of expected evaluation-period means are $[0.9947,0.9969]$ for held gross exposure and $[0.0104,0.0122]$ for absolute net exposure, in target-capital units. At least one trade hits its cap on $83.0$--$86.2\%$ of days. Mean diagnostic wealth ranges from $-1.2275$ to $-0.6305$. These are not daily quantiles or funded-account returns.\\
Rank-gradient checks & All five adjacent doublings are retained for pool B and the lowest-liquidity tercile. At 16 to 32, positive and negative local sums are $11.9492$ and $-11.8823$ in B. Their total is $0.0669$. In the tercile, the sums are $16.3527$ and $-16.3036$, with total $0.0491$. Large local terms cancel.\\
Liquidity terciles & The retained three-tercile check has positive $\Gamma_{\rm cap}$ in six of nine cells. At 500M the low-, middle-, and high-liquidity values are $0.0968,-0.0279,0.0334$. The middle tercile reverses the sign. Terciles use full-panel median liquidity. They are retrospective groups, not a test of survivor-complete data.\\
Turnover association & Primary-pool score--desired-turnover Spearman correlations range from $0.233$ to $0.307$. Across the 24 pool-by-period cells, ten correlations are negative. The turnover-product term equals $0.886$--$1.199$ times the primary 500M capacity component, with a nonzero residual. It is an algebraic attribution, not an identified cause.\\
Monotonicity and influence & The recorded gap-curve monotonicity check passes in 12 of 16 positive-capital primary cells and 29 of 72 walk-forward cells. All four primary failures are in pool A. Removing the candidate with the worst evaluation Sharpe changes the primary $\Gamma_{\rm cap}$ by at most $8.6\%$, with no sign reversal. The walk-forward relative change is below $20\%$ in 65 of 72 cells. Pool D at 25M has a net slope of $0.000665$ before removal and $0.000076$ after removal, an $88.6\%$ relative change. A small change in the gap does not imply a small relative change in a near-zero net slope.\\
\bottomrule
\end{longtable}

The execution factorial tests every combination of three participation limits, three spreads, and three impact coefficients. These give 27 settings for each pool and capital. The cost-misspecification check pairs each of the same 27 assumed settings with all 27 evaluation settings. It selects on net Sharpe at assumed capital USD 25 million and evaluates at USD 500 million. Four pools give $4\times27\times27=2{,}916$ comparisons at $N=32$.

The walk-forward selection periods start in July 2004, 2007, 2010, 2013, 2016, and 2019. Each lasts three years. Evaluation starts at the end of its paired selection period. The first five evaluation periods last three years. The last continues through the stored July 2026 endpoint. The selection windows do not overlap with one another. The evaluation windows also do not overlap with one another. However, earlier evaluation dates enter later selection windows. The period comparisons also combine changes in market conditions and liquidity. They do not isolate either cause.


The five grids start at one trial. They are $\{1,2,4,8\}$, $\{1,2,4,8,16\}$, $\{1,2,4,8,16,32\}$, $\{1,2,4,8,16,32,64\}$, and $\{1,2,4,8,16,32,64,128\}$. The 10\% rule admits the first three grids in pools A and D and all five in pools B and C. Four capital levels therefore give 64 cells. For each influence check, the candidate with the lowest evaluation Sharpe is removed. All four slope paths and rank weights are then recomputed on the reduced pool. This is a sensitivity analysis that uses evaluation outcomes, not a selection rule.

For the turnover attribution, define each candidate's capacity-related gap as uncapped gross Sharpe minus capped gross Sharpe, plus spread-only Sharpe minus net Sharpe. Regress this gap on evaluation-period mean executed turnover, with an intercept. Apply the winner weights separately to turnover and the regression residual. Linearity gives the gap slope as the regression coefficient times the selected-turnover slope, plus the selected-residual slope. The constant intercept contributes no slope. This decomposition is an identity within the recorded pool. It does not identify a causal effect of turnover.

The monotonicity check tests whether the expected gross-minus-net gap increases at every adjacent point of the six-point search grid, allowing a numerical tolerance of $10^{-15}$. It does not test whether candidate-level gaps are monotone in proxy rank.

The lowest-liquidity rank-gradient check uses pool B's configuration design with training and trading restricted to the lowest third of names by full-panel median dollar volume. The stored feature ranks remain full-universe ranks. This check does not reconstruct missing or delisted names. For the exposure diagnostics, we first calculate each candidate's mean over time. We then average those means using the exact winner probabilities. Daily exposure medians and 95th percentiles were not retained for this comparison. All sensitivity gains remain conditional on the chosen impact functions and stored price path. None validates modeled charges against observed fills.

\section{Retrospective status and limits}\label{sec:status}
The finance study has no market period that was kept entirely separate from outcome inspection. The selection-score comparisons and additional sensitivities are exploratory. Resampling specifications were recorded before the corresponding resampling outputs were inspected. Outcomes had already been inspected, and the study was not externally preregistered.

The configurations have no demonstrated trading advantage. Mean candidate evaluation Sharpe before costs is negative in all four pools. Every non-oracle selection rule in Table~\ref{tab:all-selectors} has negative mean net Sharpe. The joint capacity-component interval and the gross-slope interval include zero. More pools, capital levels, and execution settings on the same price history do not provide independent market evidence.

The survivor-filtered universe omits missing and delisted price histories. Modeled charges are not validated against observed fills. The reported series use held gross exposure for normalization and omit financing and stock-borrow charges. They therefore do not describe a funded investor account. The effect also depends on the search grid, period, and liquidity group, as the negative cases in Table~\ref{tab:finance-inventory} show.

The finite-pool rank weights apply to any fixed candidate quantity under the stated sampling rule. They do not establish the direction of an omitted-cost effect in another domain. Such a test would need a specified deployment objective, candidate-level measurements of the omitted term, and a check of its relation to proxy rank under selection-data resampling.

\bibliographystyle{plainnat}
\bibliography{refs}
\end{document}
